\chapter{Conclusion and Future Work}
\label{chap:conclusion}

This thesis asked whether a frozen, general-purpose vision-language model can be
turned into a controllable Minecraft agent cheaply --- by training only a small
action head with behavioural cloning, holding the backbone fixed and keeping the
whole ablation inexpensive enough to sweep on a single rented GPU. Using a
$2\times2$ ablation over backbone (LLaVA-1.5 versus CLIP) and language condition,
evaluated in live MineRL rollouts, it reached a clear and slightly surprising
answer: in this cheap, frozen regime the contrastive backbone is the more
effective and the more language-steerable choice \emph{under a minimal head} --- a
richer (FiLM) read-out narrows the task gap from the same LLaVA features, so the
backbone finding is best read as a head-readout-by-backbone interaction rather
than an intrinsic ranking of the two models.

Table~\ref{tab:contributions} maps each research question of
Chapter~\ref{chap:intro} to the evidence that settles it.

\begin{table}[ht]
\centering
\begin{tabular}{p{0.11\textwidth}p{0.38\textwidth}p{0.40\textwidth}}
\hline
 & Question & Finding \\
\hline
RQ1 \newline (backbone) & Used only as a frozen extractor, does a joint VLM beat
a contrastive one? & \textbf{No}, in this regime. CLIP collects $2.5$ dirt per
episode where the matched LLaVA cell is at the floor ($0.07$); the gap is a
read-out effect, since two alternative heads (FiLM, fixAB) recover CLIP-level
collection from the same LLaVA features
(Sections~\ref{sec:exp-task}, \ref{sec:exp-film}). \\
RQ2 \newline (language) & Does the prompt change behaviour, and does it depend on
the backbone? & \textbf{Yes}, and it depends entirely on the backbone. Only
CLIP + lang shifts (\texttt{attack} $-42\pm10$~pp, $h=-0.91$, in all three seeds);
both LLaVA cells and the CLIP control are flat (Section~\ref{sec:exp-cond}). The
prompt \emph{steers} behaviour but not necessarily toward the goal --- the robust
shift is a \emph{suppression} of the action chopping needs
(Section~\ref{sec:disc-chop}). \\
RQ3 \newline (evaluation) & Do rollouts reveal what frame metrics miss? &
\textbf{Yes}, decisively. The four cells are a near-tie on held-out frames yet
diverge from a dug stack of dirt to nothing in rollout
(Section~\ref{sec:exp-offline}). \\
\hline
\end{tabular}
\caption{The three research questions and the evidence that settles each.}
\label{tab:contributions}
\end{table}

Two interpretive points carry beyond the table. First, the LLaVA result is a
statement about the cheap-and-frozen regime, not about the backbone's ceiling: its
pooled text representation in fact encodes the prompt almost perfectly (a linear
probe recovers it at $99.7\%$), but in a form entangled with the image that no
head-side lever --- including a feature-gating FiLM head --- converts into
conditioning, so the bottleneck is the pooled representation and the fix is
upstream of the head (Chapter~\ref{chap:discussion}). Its visual features,
meanwhile, are usable for the task itself once the head reads them through
modulation rather than concatenation. Second, the conditioning result is a
behaviour-routing effect shown on two short-horizon tasks with the prompt doubling
as task identity; it should not be over-read as compositional
instruction-following, and absolute dirt counts mix goal-directed digging with
incidental terrain-attacking (Section~\ref{sec:exp-task}). Within those bounds the
central claim is firm: under a fixed, inexpensive training budget, a frozen
contrastive encoder with a five-million-parameter head is both a more effective
short-horizon Minecraft controller and a more language-steerable one than a frozen
seven-billion-parameter joint model \emph{read through the same minimal head}. The
task half of that gap is a head-readout effect --- a FiLM head recovers it from the
LLaVA features --- so the claim is about the cheap regime and its default
read-out, not about the backbones' intrinsic limits.

\paragraph{Outlook.} Chapter~\ref{chap:discussion} develops three directions in
detail. The first relaxes the freezing constraint on the joint backbone exactly
where this study located the bottleneck --- the entangled text representation ---
through parameter-efficient fine-tuning or a few learnable conditioning tokens,
paired with a direct representational probe to confirm whether the text features
then separate by task. The second takes the CLIP result at face value and composes
several cheap, prompt-conditioned heads into an LLM-directed hierarchy, keeping
each controller small while a planner supplies the long-horizon reasoning. The
third closes the loop with reward: the chop rollouts show a cloned CLIP agent that
already targets trees but cannot finish the block, an ideal starting point for an
online, reward-based refinement stage. The cheap-and-frozen recipe characterised
here is a natural seed for all three.
